% ID: SEL-D0001-OA-E01
% ESTADO: cotejada. Fuente y transcripción revisadas de forma independiente; salida compilada inspeccionada.
% Fuente: Selectividad/sel_2019_matematicas_aplicadas/Titular Junio/Exámenes Mat Apli CCSS A y B.pdf
% SHA-256: fa0cbbddfdf34b1f9fcf8d1090c2d0c3f13e5353e06da22ec354d3345e313584
% Página: 1; opción A; ejercicio 1.
% Encabezado original: EJERCICIO 1. Solo enunciado; no contiene solución.

\textbf{(2.5 puntos)} Una empresa textil quiere fabricar dos tipos de camisetas,
lisas y estampadas. Para fabricar una camiseta lisa necesita 70 g de algodón
y 20 g de poliéster y para cada camiseta estampada 60 g de algodón y 10 g
de poliéster. La empresa dispone para ello de 4\,200 g de algodón y 800 g
de poliéster. Para que sea rentable debe fabricar al menos 10 estampadas
y además, el doble de las estampadas debe ser al menos igual al número
de lisas. Sabiendo que cada camiseta lisa da un beneficio de 5 euros y cada
estampada de 4 euros, \textquestiondown{}cuántas camisetas de cada tipo debería
fabricar para obtener el máximo beneficio? \textquestiondown{}Cuál es ese beneficio?
